\section{Localized modes and self-adjoint eigenproblems}
\label{sec:localized}

A core goal of the substrate theory is to explain particle-like behavior (localization, discrete spectra,
internal ``spin/charge''-like structure) as properties of confined wave states rather than point particles.

\subsection{Linearization about symmetric backgrounds and self-adjointness}

Let $\vecX_0(x)$ be a time-independent background embedding representing a localized deformation region.
We are especially interested in backgrounds that are \emph{spherically symmetric} in the intrinsic brane coordinates,
so that invariants depend only on $r=|x|$.

Linearize the equations of motion \eqref{eq:eom} about $\vecX_0$ by writing
$\vecX(x,t)=\vecX_0(x)+\bm{\eta}(x,t)$ with small perturbation $\bm{\eta}$.
Because the dynamics derives from an action \eqref{eq:action}, the linearized operator arises from the second
variation of the action and is therefore self-adjoint with respect to the inner product induced by the kinetic energy,
schematically
\begin{equation}
  \langle \bm{\eta},\bm{\zeta}\rangle
  :=\int_{\Omega} \rho_m\,\bm{\eta}(x)\cdot \bm{\zeta}(x)\,\dd^3x,
  \label{eq:inner-product}
\end{equation}
together with the boundary/regularity conditions appropriate to the chosen domain (periodic cell, decay at infinity,
or a large but finite domain with physical boundary conditions).
In this setting, orthogonality and completeness of eigenmodes are not ``imposed'' as ad-hoc constraints: they follow as
standard properties of a self-adjoint eigenproblem.

\subsection{Angular structure and spherical harmonics}

For a spherically symmetric background, the linearized operator commutes with the action of $SO(3)$ on the intrinsic
coordinates. Consequently, eigenmodes can be chosen to transform in irreducible $SO(3)$ representations.
Operationally, this implies a separation into a radial problem and an angular problem whose eigenfunctions are
spherical harmonics $Y_{\ell m}(\theta,\varphi)$.
This statement is a symmetry consequence and does not depend on switching the substrate constitutive law:
it only requires rotational invariance of the background and of the linearized operator.

In particular, spherical-harmonic mode families provide the natural spatial ``skeleton'' for localized excitations.
Appropriate superpositions (and/or coupling to embedding-direction components) can yield torus-like or ring-like
energy density profiles reminiscent of toroidal electron models, but here without assuming literal multi-winding
of a photon trajectory.

\subsection{Amplitude scaling and energy normalization}

The linear eigenproblem fixes mode shapes and frequencies but not the overall amplitude:
if $\bm{\eta}$ is an eigenmode, so is $c\,\bm{\eta}$ for any scalar $c$.
Physical amplitudes are selected only when nonlinearities are included and the total energy is fixed.
For an electron-like excitation one may impose a rest-energy normalization
\begin{equation}
  \overline{E}[\vecX] \;=\; m_e c^2,
  \label{eq:rest-energy-normalization}
\end{equation}
where $\overline{E}$ denotes the conserved substrate energy associated with \eqref{eq:action} and $c$ is the
characteristic speed of the carrier branch used to define the effective spacetime coordinates in \Cref{sec:geophase}.
This ties the nonlinear amplitude scale of a localized soliton to a physical energy target.

\subsection{Confinement as nonlinear self-guidance}

A confining mechanism is required: in a purely linear wave system, localized packets generally disperse or radiate.
A useful conceptual benchmark is the toroidal electron proposal attributed to Williamson \& van der Mark, which argues
that ordinary gravitational curvature is far too weak to confine a Compton-scale photon-like mode and that confinement
must arise from a self-consistent, excitation-specific mechanism.
We adopt the \emph{confinement requirement} while not committing to a literal double-loop trajectory
(see e.g.\ the review discussion in \citep{HuygensOptics2025}).

In the present substrate theory, confinement is attributed primarily to the brane's \emph{constitutive and geometric
nonlinearity}: large-amplitude deformations modify the local stress response and therefore the local linearized
propagation characteristics of small fluctuations around the deformed state.
Equivalently, the soliton ``writes'' an effective medium for itself (through $g_{ij}(\vecX)$ and $W$), enabling
self-guidance and trapping without introducing a fundamental point charge or external potential.

\subsection{Connection to Berry/Wilczek--Zee holonomy and spin}

Localized modes can carry internal polarization structure (multi-component standing waves).
The Berry/Wilczek--Zee connections of \Cref{sec:geophase} can then be evaluated on loops surrounding the localized region
or on adiabatic deformations of a localized mode family.
In this way, long-range ``charge''-like labels are interpreted as far-field holonomy/phase structure, while the
spin-$\tfrac12$ $4\pi$ aspect is attributed to spinorial holonomy of a two-level polarization subspace.
